\documentclass[11pt]{article}

\usepackage[T1]{fontenc}
\usepackage[utf8]{inputenc}
\usepackage{amsmath,amssymb,amsthm}
\usepackage{array,booktabs}
\usepackage[margin=1in]{geometry}
\usepackage[hidelinks]{hyperref}

\newtheorem{theorem}{Theorem}
\newtheorem{lemma}[theorem]{Lemma}
\newtheorem{proposition}[theorem]{Proposition}
\newtheorem{corollary}[theorem]{Corollary}
\theoremstyle{definition}
\newtheorem{definition}[theorem]{Definition}
\newtheorem{remark}[theorem]{Remark}

\newcommand{\TS}{\mathsf{TS}}
\newcommand{\BPTS}{\mathsf{BPTS}}
\newcommand{\BPTIME}{\mathsf{BPTIME}}
\newcommand{\MATIME}{\mathsf{MATIME}}
\newcommand{\QCMATIME}{\mathsf{QCMATIME}}
\newcommand{\leanus}{\textunderscore\allowbreak}
\newcommand{\lean}[1]{\ifmmode\text{\ttfamily #1}\else{\ttfamily\let\_\leanus #1}\fi}
\newcommand{\pin}{\mathrm{pin}}
\newcommand{\pot}{\mathrm{pot}}
\newcommand{\up}{U}

\title{The Mudigonda--Williams calculus for $\exists\cdot\BPTIME$ versus $\BPTS$\\ stops exactly at the golden ratio: a kernel-checked threshold}
\author{PastaPastaPasta}
\date{November 2026}

\emergencystretch=2em

\begin{document}
\maketitle

\begin{abstract}
Mudigonda and Williams \cite{MW21} define a calculus of alternation-trading proofs for showing that problems verifiable by Merlin--Arthur protocols in linear time cannot be solved by randomized algorithms running in time $n^c$ and space $n^{o(1)}$. Their Theorem~5.3 proves, by a proof in this calculus, that $\exists\cdot\BPTIME[n]\not\subseteq\BPTS[n^c]$ for every $c<r_1\approx1.4656$, the largest root of $x^3-x^2-1$.

We determine the exact threshold of a Lean~4 formalisation of this calculus. For every real $c\ge1$, a refutation exists if and only if $c(c-1)<1$, that is, if and only if $c<\varphi=(1+\sqrt5)/2\approx1.6180$. The construction for $c<\varphi$ repeats one move of \cite{MW21}: a randomized speedup applied under the outer quantifier block that the first randomized speedup leaves behind. The source paper makes this move itself in the proof of a later theorem, its quantum result (Theorem~5.4). The impossibility for $c\ge\varphi$ is a potential-function argument on a relaxed calculus. The relaxed calculus also contains the unmerged orientation of the randomized speedup, speedups at the end points of their parameter range, and weakening of exponents.

\emph{What is kernel-checked:} the threshold of the \emph{formal} calculus, in both directions, including a 15-step refutation at $c=1.467>r_1$.
\emph{What is not machine-checked:} that the formal rules coincide with the rules of \cite{MW21} (a modelling claim, which we compare rule by rule on paper in Table~\ref{tab:rules}), and the soundness of the lemmas of \cite{MW21}, on which the published bound $1.4656$ depends in the same way. This includes its hierarchy and contradiction lemma (Lemma~5.2), whose proof \cite{MW21} only sketch. Conditional on exactly the lemmas that \cite[Theorem~5.3]{MW21} uses, the formal result gives $\exists\cdot\BPTIME[n]\not\subseteq\BPTS[n^c]$, hence $\MATIME[n]\not\subseteq\BPTS[n^c]$, for every $c<\varphi$. The threshold is a limit on one proof calculus. It says nothing about lower bounds proved by other methods.
\end{abstract}

\section{Introduction}

\paragraph{Background.}
Alternation-trading proofs \cite{W13} derive a contradiction from a hypothesis such as ``SAT has an algorithm running in time $n^c$ and space $n^{o(1)}$'' by alternately applying \emph{speedup} rules, which add quantifiers to a class and reduce its running time, and \emph{slowdown} rules, which use the hypothesis to remove a quantifier at the cost of a factor $c$ in the exponent. Fortnow and van Melkebeek obtained the exponent $\varphi$ for SAT against deterministic algorithms in this way \cite{FvM00,FLvMV05}; Williams later reached $2\cos(\pi/7)\approx1.8019$, and Buss and Williams \cite{BW15} proved that $2\cos(\pi/7)$ is the limit of alternation-trading proofs with those rules.

Mudigonda and Williams \cite{MW21} extend the method to randomized and quantum verifiers. For randomized algorithms, the hypothesis is $H_c$: $\exists\cdot\BPTIME[n]\subseteq\BPTS[n^c]$, where $\BPTS[t]$ is two-sided-error randomized time $t^{1+o(1)}$ with space $n^{o(1)}$ (the ITCS version uses $O(\log n)$ space; nothing below depends on the difference). Their speedups are the Diehl--van Melkebeek randomized speedup \cite{DvM06} (their Theorem~2.10 and Corollary~2.11) and the usual deterministic speedup (Corollary~2.8). The slowdown is their Lemma~5.1, and their Lemma~5.2 turns a proof $\BPTS[n^d]\subseteq\dots\subseteq\BPTS[n^{d'}]$ with $d'<d$ into a contradiction. Their Theorem~5.3 gives $\exists\cdot\BPTIME[n]\not\subseteq\BPTS[n^c]$ for $c<r_1\approx1.4656$, the largest root of $x^3-x^2-1$, and hence $\MATIME[n]\not\subseteq\BPTS[n^c]$ (part of their Main Theorem~3). The previous bound was $\sqrt2$ \cite{DvM06,Die07}. \cite{MW21} make no optimality claim for $r_1$; their optimality result (Appendix~A) concerns their Main Theorem~1, for deterministic verifiers.

Result numbers of the form $x.y$ (Theorem~2.10, Lemma~5.2, \dots) refer to the arXiv version of \cite{MW21}.\footnote{In the ITCS version: Definition~2.14 is Definition~12, Lemma~2.7 is Lemma~4, Corollary~2.8 is Corollary~5, Theorem~2.10 is Theorem~7, Corollary~2.11 is Corollary~8, Lemma~2.17 is Lemma~15, Lemmas~5.1 and~5.2 are Lemmas~33 and~34, Theorems~5.3 and~5.4 are Theorems~35 and~36.} Our own results are numbered $1,2,3,\dots$ All theorems below are Lean theorems; names refer to the accompanying Lean library (namespace \lean{PastaMath.MWThreshold}).

\paragraph{Results.}
\begin{itemize}
\item \textbf{Threshold} (Theorem~\ref{thm:threshold}; \lean{MWG.threshold\_phi}, \lean{MWG.threshold\_goldenRatio}). For $c\ge1$, the formal calculus \lean{MWL.Step} (Section~\ref{sec:calculus}) has a refutation if and only if $c(c-1)<1$, i.e.\ $c<\varphi$.
\item \textbf{Construction} (Theorem~\ref{thm:construction}; \lean{MWG.refutes\_of\_lt\_phi}). For $1\le c<\varphi$ a refutation exists in the exact rules. A refutation with 15 steps at $c=1.467>r_1$ is checked separately (\lean{MWG.mw15\_ok}, \lean{MWG.mw21\_root\_lt}; Section~\ref{sec:fifteen}).
\item \textbf{Limit} (Theorem~\ref{thm:limit}; \lean{MW.no\_mrefutation}, \lean{MW.no\_urefutation}, \lean{MW.no\_wrefutation}, \lean{MWL.no\_lit\_refutation}). For $c\ge\varphi$ there is no refutation, strict or not, in a relaxed calculus that contains the exact rules, the unmerged orientation of the randomized speedup, and weakening. Every quantifier-free line reached from $\BPTS[n^{a_0}]$ has exponent at least $c\,a_0$.
\end{itemize}
Conditional on the lemmas of \cite{MW21} listed in Section~\ref{sec:consequence}, this gives $\exists\cdot\BPTIME[n]\not\subseteq\BPTS[n^c]$ for every $c<\varphi$ (Corollary~\ref{cor:lb}).

\paragraph{Where the improvement comes from.}
The proof of \cite[Theorem~5.3]{MW21} uses one randomized speedup (Theorem~2.10), then $k-1$ deterministic speedups and $k+2$ slowdowns. The first randomized speedup leaves an outer block $(\exists\,n^0)^1$. Before the last slowdown removes it, the line is $(\exists\,n^0)^1\BPTS[n^{e}]$, and the randomized speedup Corollary~2.11 applies to it again; its own leading $n^0$ block is absorbed into $(\exists\,n^0)^1$ (the footnote to Corollary~2.11). Repeating the inner pattern under this block contracts the exponent by a factor $\rho<1$ whenever $c(c-1)<1$ (for a long enough pattern), and the factor $c$ for removing the block is paid only once. Using the pattern once needs $c\cdot c(c-1)<1$ in the limit, which is $c^3-c^2-1<0$. \cite{MW21} repeat a pattern under a persistent outer block in the proof of their Theorem~5.4 (for $\exists\cdot\mathsf{BQTIME}$, with Grover search in place of the classical speedups); the construction here is the same move in the classical setting.

\paragraph{Novelty.}
We are not aware of a prior statement of this improvement, or of a determination of the threshold of this calculus. Our search was not exhaustive.

\paragraph{AI assistance.}
The proofs and the Lean formalisation were produced with Anthropic's Claude models under the author's direction. The kernel-checked Lean proofs are the primary evidence for the theorems about the formal calculus; this text explains them and states what they do and do not establish.

\section{The calculus}\label{sec:calculus}

\subsection{Lines}

A \emph{line} is an alternating complexity class
\[
(Q_1\,n^{a_1})^{b_1}\,(Q_2\,n^{a_2})^{b_2}\cdots(Q_k\,n^{a_k})^{b_k}\;K[n^{d}]
\]
\cite[Definition~2.5]{MW21}: in stage $i$ the machine receives a string of length $n^{a_i}$ from quantifier $Q_i\in\{\exists,\forall\}$ and passes $n^{b_i}$ bits to the next stage; the verifier $K\in\{\TS,\BPTS\}$ runs in time $n^{d}$; all $n^{o(1)}$ factors are suppressed. We call $a_i$ the \emph{guess exponent} and $b_i$ the \emph{output exponent} (or superscript) of block $i$, and set $b_0=1$ (the input has length $n$). A block written $(Q\,x\log n)^b$ or $(Q\,n^0)^b$ has guess exponent $0$. For a block list $P$ let $\pin(P)$ be the output exponent of its innermost block, with $\pin(\emptyset)=1$. Every line of a proof is \emph{orderly}, $a_i\le b_i$ \cite[Lemma~2.17]{MW21}.

\subsection{Rules}\label{sec:rules}

Fix $c$. The rules are those of \cite[Definition~2.14]{MW21} that can occur in a proof from a $\BPTS$ line, with the slowdown of \cite[Lemma~5.1]{MW21} (the generic slowdown of Definition~2.13 with $\alpha=1$). In each rule $P$ is the rest of the block list and $(Q\,n^{a})^{b}$ the innermost block.

\begin{enumerate}
\item[(R3)] \textbf{Randomized speedup on a quantifier-free line} (rule~3, Theorem~2.10). For $0<x<d$:
\[ \BPTS[n^d]\subseteq(\exists\,n^0)^1(\forall\,n^x)^{\max(x,1)}(\exists\,x\log n)^1\,\TS[n^{d-x}]. \]
\item[(R4)] \textbf{Randomized speedup under a block} (rule~4, Corollary~2.11). For $0<x<d$:
\[ P\,(Q\,n^{a})^{b}\,\BPTS[n^d]\subseteq P\,(Q\,n^{a})^{b}(\neg Q\,n^{x})^{\max(b,x)}(Q\,x\log n)^{b}\,\TS[n^{d-x}]. \]
Corollary~2.11 prints the new quantifiers as $Q_{k+1},Q_{k+2}$. By its footnote, the leading $n^0$ quantifier of Theorem~2.10 ``is absorbed into the previous quantifier'', which forces $Q_{k+1}=\neg Q$ and $Q_{k+2}=Q$.
\item[(R2)] \textbf{Deterministic speedup} (rule~2, Corollary~2.8). For $0<x<d$:
\[ P\,(Q\,n^{a})^{b}\,\TS[n^d]\subseteq P\,(Q\,n^{\max(a,x)})^{\max(b,x)}(\neg Q\,x\log n)^{b}\,\TS[n^{d-x}]. \]
\item[(S)] \textbf{Slowdown} (rule~5, Lemma~5.1, under $H_c$). For a kernel $K\in\{\TS,\BPTS\}$ (a $\TS$ kernel is a $\BPTS$ kernel):
\[ P\,(Q\,n^{a})^{b}\,K[n^d]\subseteq P\,\BPTS[n^{c\max(d,\,b,\,a,\,\pin(P))}]. \]
\end{enumerate}
Rule~1 (Lemma~2.7 on a quantifier-free $\TS$ line) never applies: the start line is $\BPTS$, a $\TS$ kernel is produced only by a speedup, and every slowdown produces a $\BPTS$ kernel. For the same reason a deterministic speedup is applied only to a line produced by a speedup (randomized or deterministic), never directly after a slowdown, and a randomized speedup only to a line produced by a slowdown or to the start line.

\begin{definition}[refutation; \lean{MWL.Refutes}]\label{def:refutes}
A \emph{refutation} at $c$ is a sequence of $N\ge1$ rule applications from $\BPTS[n^{a_0}]$, $a_0>0$, to a quantifier-free line $\BPTS[n^{a'}]$ with $a'<a_0$.
\end{definition}
Intermediate lines are unrestricted. The first step is necessarily (R3), so a refutation has at least two inclusions, as Lemma~5.2 requires. Lemma~5.2 as printed says ``$d<d'$''; Definition~2.14 says $d'\le d$, and the proof of Lemma~5.2 uses $d'<d$. We read the condition as strict; the limit theorem holds for the non-strict version as well.

\subsection{The formalisation}

A Lean line is a triple $(\mathit{bl},k,a)$: the block list $\mathit{bl}$, stored innermost first, of blocks $\langle q,g,o\rangle$ ($q=\mathtt{true}$ for $\exists$, $g$ the guess exponent, $o$ the output exponent); the kernel type $k$ ($\mathtt{true}$ for $\BPTS$); and the kernel exponent $a\in\mathbb R$ (\lean{MWL.Line}, \lean{MWL.Blk}, \lean{MWL.pin}). The relation \lean{MWL.Step}~$c$ has one constructor per rule, with the exact output values and side conditions above; \lean{MWL.Refutes} is Definition~\ref{def:refutes}. Table~\ref{tab:rules} compares them with \cite{MW21}.

\begin{table}[ht]
\centering
\small
\begin{tabular}{@{}>{\raggedright\arraybackslash}p{3.3cm}>{\raggedright\arraybackslash}p{2.7cm}>{\raggedright\arraybackslash}p{8.5cm}@{}}
\toprule
\cite{MW21} & Lean & Effect, cost and side conditions in \lean{MWL.Step} ($u=\langle q,g,o\rangle$ innermost block, $\mathit{rest}$ the blocks outside it) \\
\midrule
Def.~2.14 rule~3, Theorem~2.10 & \lean{rspeed0} &
$([],\BPTS,a)\mapsto([\langle\exists,0,1\rangle,\langle\forall,x,\max(x,1)\rangle,\langle\exists,0,1\rangle],\TS,a-x)$ (innermost first); requires $0<x<a$ \\
Def.~2.14 rule~4, Corollary~2.11 (merged, footnote~6) & \lean{rspeed} &
$(u::\mathit{rest},\BPTS,a)\mapsto(\langle q,0,o\rangle::\langle\neg q,x,\max(o,x)\rangle::u::\mathit{rest},\TS,a-x)$; $u$ unchanged; requires $0<x<a$ \\
Def.~2.14 rule~2, Corollary~2.8 & \lean{dspeed} &
$(u::\mathit{rest},\TS,a)\mapsto(\langle\neg q,0,o\rangle::\langle q,\max(g,x),\max(o,x)\rangle::\mathit{rest},\TS,a-x)$; requires $0<x<a$ \\
Def.~2.14 rule~5, Lemma~5.1 ($\alpha=1$) & \lean{slow} &
$(u::\mathit{rest},k,a)\mapsto(\mathit{rest},\BPTS,c\max(a,o,g,\pin(\mathit{rest})))$ for either kernel; $\pin([])=1$ \\
Lemma~5.2, Def.~2.14 & \lean{Refutes} &
$N\ge1$ steps from $([],\BPTS,a_0)$, $a_0>0$, to $([],\BPTS,a')$ with $a'<a_0$ \\
\midrule
Theorem~2.10 under a block, leading $n^0$ block not absorbed & \lean{StepU.rspeedU} &
$(u::\mathit{rest},\BPTS,a)\mapsto(\langle\neg q,0,o\rangle::\langle q,x,\max(o,x)\rangle::\langle\neg q,0,o\rangle::u::\mathit{rest},\TS,a-x)$; requires $0<x<a$ \\
trivial inclusions (weakening) & \lean{StepW.weaken} &
same quantifiers and block count; any guess exponent, output exponent or the kernel exponent raised (guess $\le$ output); a $\TS$ kernel may become $\BPTS$ \\
\bottomrule
\end{tabular}
\caption{The rules of \cite{MW21} used in proofs from a $\BPTS$ line, and their Lean counterparts. Constructor names are those of \lean{MWL.Step}, \lean{MWL.StepU} and \lean{MWL.StepW}. Rows below the middle line are not rules of \cite{MW21}; they are covered by the limit theorem only (Section~\ref{sec:coverage}).}
\label{tab:rules}
\end{table}

\paragraph{Reading points.}
The following choices are ours; each is either forced by the text of \cite{MW21} or affects only the limit side, where it is covered anyway (Section~\ref{sec:coverage}).
\begin{enumerate}
\item \emph{Rule~3's last block.} \cite{MW21} do not print the output exponent of the $(\exists\,x\log n)$ block of rule~3. We use $1$, the value $b_k$ that Corollary~2.11 gives for $k=0$ (the input). The relaxed calculus allows any value.
\item \emph{Orientation of rule~4.} We use the merged orientation prescribed by the footnote to Corollary~2.11. The construction uses only it; the limit also covers the other one.
\item \emph{Theorem~2.10's orientation.} Theorem~2.10 is printed as $\BPTS\subseteq(\forall\,n^0)^1(\exists\,n^x)(\forall\,x\log n)\TS$, rule~3 of Definition~2.14 as $(\exists\,n^0)^1(\forall\,n^x)(\exists\,x\log n)\TS$; they are equivalent because $\BPTS$ is closed under complement.
\item \emph{Slowdown of a $\TS$ kernel.} Lemma~5.1 is stated for $\BPTS$ kernels; rule~5 of Definition~2.14 is stated for $(\mathsf{BP})\TS$ kernels, and \cite[Theorem~5.3]{MW21} applies it to a $\TS$ kernel. We use $\TS\subseteq\BPTS$.
\item \emph{Strictness} (Definition~\ref{def:refutes}).
\item \emph{Lower-order terms.} As in \cite{MW21}, $n^{o(1)}$ factors are suppressed and do not appear in the exponents.
\end{enumerate}

\section{Statements}

\begin{theorem}[\lean{MWG.threshold\_phi}, \lean{MWG.threshold\_goldenRatio}]\label{thm:threshold}
For every real $c\ge1$:
\[ \bigl(\exists N,M:\ \lean{MWL.Refutes}\ c\ N\ M\bigr)\iff c(c-1)<1\iff c<\tfrac{1+\sqrt5}{2}. \]
\end{theorem}

\begin{theorem}[\lean{MWG.refutes\_of\_lt\_phi}]\label{thm:construction}
If $c\ge1$ and $c(c-1)<1$, there is a refutation at $c$ in \lean{MWL.Step}.
\end{theorem}

\begin{theorem}[\lean{MW.qf\_ge}, \lean{MW.no\_mrefutation}, \lean{MWL.no\_lit\_refutation}; \lean{MW.no\_urefutation}, \lean{MW.no\_wrefutation}, \lean{MWL.no\_litU\_refutation}, \lean{MWL.no\_litW\_refutation}]\label{thm:limit}
Let $c>0$ with $c^2\ge c+1$. In the relaxed calculus of Section~\ref{sec:limit}, every quantifier-free line $\BPTS[n^{a'}]$ reached from $\BPTS[n^{a_0}]$, $a_0\ge0$, by at least one rule step (and any number of weakening steps) has $a'\ge c\,a_0$. Hence there is no refutation, strict or non-strict, in \lean{MWL.Step}, in its extension \lean{MWL.StepU} by the unmerged orientation of rule~4, or in its further extension \lean{MWL.StepW} by weakening.
\end{theorem}

Theorem~\ref{thm:threshold} is Theorems~\ref{thm:construction} and~\ref{thm:limit} together; $c(c-1)<1$ is equivalent to $c<\varphi$ for $c\ge1$. The same equivalence holds with \lean{StepU} and with \lean{StepW} in place of \lean{Step} (\lean{MWG.threshold\_phi\_U}, \lean{MWG.threshold\_phi\_W}; for \lean{StepW} the refutation is non-strict and must contain at least one rule step, since a weakening alone gives $\BPTS[n^{a_0}]\subseteq\BPTS[n^{a_0}]$).

\section{The construction for \texorpdfstring{$c<\varphi$}{c < phi}}\label{sec:construction}

Fix $c\ge1$. Let $\sigma_0=1$ and $\sigma_{i+1}=\sigma_i+c^i$, so $\sigma_j=1+\sum_{i<j}c^i$ and $(c-1)\sigma_j=c^j+c-2$ (\lean{MWG.sig}, \lean{MWG.sig\_closed}). Write $\top=(\exists\,n^0)^1$ for the outer block that (R3) creates (\lean{MWG.top}).

\begin{lemma}[inner phase; \lean{MWG.inner}]\label{lem:inner}
Let $j\ge0$, $t\ge1$, and let $R$ be a block list with $\pin(R)\le c^jt$. Then the line $R\,(q\,\log n)^1\,\TS[n^{\sigma_jt}]$ derives $R\,\BPTS[n^{v}]$ for some $0\le v\le c^{j+1}t$, using $j$ deterministic speedups and $j+1$ slowdowns.
\end{lemma}
\begin{proof}
Induction on $j$. For $j=0$ the kernel is $t$; one slowdown costs $c\max(t,1,0,\pin(R))\le ct$, since $t\ge1$ and $\pin(R)\le t$. For $j+1$ the kernel is $\sigma_{j+1}t=\sigma_jt+c^jt$. Apply (R2) with $x=c^jt$, which satisfies $1\le x<\sigma_{j+1}t$: the block $(q\,\log n)^1$ becomes $(q\,n^{x})^{x}$ (as $\max(0,x)=x=\max(1,x)$), a new block $(\neg q\,\log n)^1$ is opened, and the kernel becomes $\sigma_jt$. By induction with $R'=R\,(q\,n^x)^x$, $\pin(R')=c^jt$, we reach $R\,(q\,n^x)^x\,\BPTS[n^{v}]$ with $v\le c^{j+1}t$. A slowdown costs $c\max(v,x,x,\pin(R))\le c\cdot c^{j+1}t$, because $x=c^jt\le c^{j+1}t$ and $\pin(R)\le c^{j+1}t$.
\end{proof}

All the floors $\max(\cdot,1)$ in the rules are resolved by $t\ge1$: every speedup parameter is at least $t\ge1$, and every kernel exponent in the derivation is at least $t$.

\begin{lemma}[the macro $M(j+1)$; \lean{MWG.mac}]\label{lem:mac}
Let $j\ge0$, $T\ge1$ and $e=\sigma_{j+1}T$. Both $\BPTS[n^{e}]$ and $\top\,\BPTS[n^{e}]$ derive $\top\,\BPTS[n^{e'}]$ for some $0\le e'\le c^{j+2}T=\rho_{j+1}\,e$, where
\[ \rho_{j+1}=\frac{c^{j+2}}{\sigma_{j+1}}=\frac{c^{j+2}(c-1)}{c^{j+1}+c-2}\quad(c\neq1). \]
The derivation (the \emph{macro} $M(j+1)$) is one randomized speedup with $x=c^jT$, then the deterministic speedups $x=c^{j-1}T,\dots,cT,T$, then $j+2$ slowdowns.
\end{lemma}
\begin{proof}
From $\BPTS[n^e]$ apply (R3), from $\top\,\BPTS[n^e]$ apply (R4) under $\top$ (its new $n^0$ block is absorbed into $\top$), with $x=c^jT\in[1,e)$. Both give $\top\,(\forall\,n^{x})^{x}(\exists\,\log n)^1\,\TS[n^{\sigma_jT}]$. Lemma~\ref{lem:inner} with $R=\top\,(\forall\,n^x)^x$, $\pin(R)=c^jT$, reaches $\top\,(\forall\,n^x)^x\,\BPTS[n^v]$ with $v\le c^{j+1}T$, and the last slowdown costs $c\max(v,x,x,1)\le c^{j+2}T$, because $1\le x\le c^{j+1}T$.
\end{proof}

The kernel after all speedups of the macro is exactly $T$, and the slowdowns cost at most $cT,c^2T,\dots,c^{j+2}T$ in turn.

\paragraph{Example: $c=3/2$.}
Here $\sigma_2=1+1+\tfrac32=\tfrac72$ and $\rho_2=(3/2)^3/(7/2)=27/28$. For $e\ge7/2$, $T=2e/7$ and the macro $M(2)$ is
\begin{align*}
\top\,\BPTS[n^{e}]
&\subseteq\top\,(\forall\,n^{3e/7})^{3e/7}(\exists\,\log n)^1\,\TS[n^{4e/7}] &&\text{Cor.~2.11, }x=3e/7\\
&\subseteq\top\,(\forall\,n^{3e/7})^{3e/7}(\exists\,n^{2e/7})^{2e/7}(\forall\,\log n)^1\,\TS[n^{2e/7}] &&\text{Cor.~2.8, }x=2e/7\\
&\subseteq\top\,(\forall\,n^{3e/7})^{3e/7}(\exists\,n^{2e/7})^{2e/7}\,\BPTS[n^{3e/7}] &&\text{Lemma~5.1: }\tfrac32\max(\tfrac{2e}7,1,0,\tfrac{2e}7)\\
&\subseteq\top\,(\forall\,n^{3e/7})^{3e/7}\,\BPTS[n^{9e/14}] &&\text{Lemma~5.1: }\tfrac32\max(\tfrac{3e}7,\tfrac{2e}7,\tfrac{2e}7,\tfrac{3e}7)\\
&\subseteq\top\,\BPTS[n^{27e/28}] &&\text{Lemma~5.1: }\tfrac32\max(\tfrac{9e}{14},\tfrac{3e}7,\tfrac{3e}7,1).
\end{align*}
Each application multiplies the kernel exponent by exactly $27/28$ (the script \texttt{scripts/mw\_check.py macro} checks this exactly). Twelve applications from $a_0=25/4$ (the first one creates $\top$ by (R3)) and the final slowdown give a 61-step refutation, ending at $\BPTS[n^{6.0595\ldots}]$.

\begin{lemma}[\lean{MWG.exists\_contracting}, via Bernoulli's inequality \lean{MWG.bernoulli}]\label{lem:contract}
If $c\ge1$ and $c(c-1)<1$, then $\rho_{j+1}<1$, i.e.\ $c^{j+2}<\sigma_{j+1}$, for some $j$.
\end{lemma}
\begin{proof}
For $c=1$ take $j=0$: $\rho_1=1/2$. For $c>1$, using $(c-1)\sigma_{N+1}=c^{N+1}+c-2$, the claim $c^{N+2}<\sigma_{N+1}$ is equivalent to $c^{N+1}(1+c-c^2)>2-c$. Here $1+c-c^2>0$, and $c^{N+1}\ge1+(N+1)(c-1)$ grows without bound.
\end{proof}

As $j\to\infty$, $\rho_{j+1}\to c(c-1)$; so the macros contract for some length exactly when $c(c-1)<1$. The shortest contracting macro at $c$ is $M(1)$ (with $\rho_1=c^2/2$) for $c<\sqrt2$, $M(2)$ (with $\rho_2=c^3/(c+2)$) for $c$ below the real root $1.5213\ldots$ of $c^3=c+2$, and longer macros closer to $\varphi$. Iterating $M(1)$ recovers the earlier bound $\sqrt2$.

\begin{proof}[Proof of Theorem~\ref{thm:construction}]
Take $j$ from Lemma~\ref{lem:contract}, write $\sigma=\sigma_{j+1}$ and $\rho=\rho_{j+1}<1$, and choose $m$ with $c\rho^{m+1}<1$. Start from $a_0=c\sigma+1$. One macro from $\BPTS[n^{a_0}]$ (with $T=a_0/\sigma\ge1$) reaches $\top\,\BPTS[n^{e_1}]$ with $e_1\le\rho a_0$. Repeat the macro under $\top$ while the kernel $e$ satisfies $e\ge\sigma$ (so $T=e/\sigma\ge1$), at most $m$ times; each round multiplies the kernel by at most $\rho$, so afterwards $e'\le\max(\rho^{m}e_1,\sigma)\le\max(\rho^{m+1}a_0,\sigma)$ (\lean{MWG.iter}). The final slowdown removes $\top$ at cost $c\max(e',1,0,1)\le c\max(\rho^{m+1}a_0,\sigma)<a_0$, since $c\rho^{m+1}<1$ and $c\sigma<a_0$.
\end{proof}

\paragraph{Side conditions.}
Every speedup parameter in the construction is at least $1$ and below the current kernel exponent, every kernel exponent is at least $1$, and every line is orderly and alternates. The outer block $\top$ is never popped until the last step. Its guess exponent enters only the final slowdown; the slowdowns that remove the block directly inside $\top$ use its output exponent $1$.

\paragraph{Relation to the proof of \cite[Theorem~5.3]{MW21}.}
The annotation $1^k0^{k+2}$ of \cite[Theorem~5.3]{MW21} is the macro $M(k)$ used once from $\BPTS[n^d]$, followed by the final slowdown. It is a refutation when $c^{k+2}<\sigma_k$, i.e.\ $c\rho_k<1$, and in the limit $k\to\infty$ this is $c^2(c-1)<1$, i.e.\ $c<r_1$. The parameters printed in \cite{MW21} differ from these (Section~\ref{sec:errata}). The proof of \cite[Theorem~5.3]{MW21} never applies Corollary~2.11. The construction above applies it every time it repeats the macro under $\top$; the first use is the eighth line of the 15-step proof below.

\subsection{A 15-step refutation at \texorpdfstring{$c=1.467$}{c = 1.467}}\label{sec:fifteen}

The sequence $M(3)\,M(3)$ followed by the final slowdown is a refutation at $c=1467/1000$ from $a_0=10^4$ (\lean{MWG.mw15\_ok}, checked by kernel evaluation of the executable checker \lean{MWL.checkQ}, whose soundness is \lean{MWL.checkQ\_sound}). Every real root of $x^3-x^2-1$ is below $1.467$ (\lean{MWG.mw21\_root\_lt}), so this proof works at an exponent where \cite[Theorem~5.3]{MW21} does not apply. The speedup parameters are not those of the general construction; they were chosen for this sequence, and any parameters satisfying the side conditions may be used. Exponents are rounded to two decimals; the checks are exact. Lines are outermost first.

{\footnotesize
\begin{align*}
 0\quad&\BPTS[n^{10000}]\\
 1\quad&\subseteq(\exists\,n^0)^1(\forall\,n^{3830})^{3830}(\exists\,\log n)^1\,\TS[n^{6170}] &&\text{Thm~2.10, }x=3830\\
 2\quad&\subseteq\cdots(\forall\,n^{3830})^{3830}(\exists\,n^{2610.76})^{2610.76}(\forall\,\log n)^1\,\TS[n^{3559.24}] &&\text{Cor.~2.8, }x=206250/79\\
 3\quad&\subseteq\cdots(\exists\,n^{2610.76})^{2610.76}(\forall\,n^{1779.66})^{1779.66}(\exists\,\log n)^1\,\TS[n^{1779.58}] &&\text{Cor.~2.8, }x=105000/59\\
 4\quad&\subseteq(\exists\,n^0)^1(\forall\,n^{3830})^{3830}(\exists\,n^{2610.76})^{2610.76}(\forall\,n^{1779.66})^{1779.66}\,\BPTS[n^{2610.76}] &&\text{Lemma~5.1}\\
 5\quad&\subseteq(\exists\,n^0)^1(\forall\,n^{3830})^{3830}(\exists\,n^{2610.76})^{2610.76}\,\BPTS[n^{3829.99}] &&\text{Lemma~5.1}\\
 6\quad&\subseteq(\exists\,n^0)^1(\forall\,n^{3830})^{3830}\,\BPTS[n^{5618.61}] &&\text{Lemma~5.1}\\
 7\quad&\subseteq(\exists\,n^0)^1\,\BPTS[n^{8242.50}] &&\text{Lemma~5.1}\\
 8\quad&\subseteq(\exists\,n^0)^1(\forall\,n^{3156.82})^{3156.82}(\exists\,\log n)^1\,\TS[n^{5085.68}] &&\text{Cor.~2.11, }x=1550000/491\\
 9\quad&\subseteq\cdots(\forall\,n^{3156.82})^{3156.82}(\exists\,n^{2151.90})^{2151.90}(\forall\,\log n)^1\,\TS[n^{2933.78}] &&\text{Cor.~2.8, }x=170000/79\\
10\quad&\subseteq\cdots(\exists\,n^{2151.90})^{2151.90}(\forall\,n^{1466.85})^{1466.85}(\exists\,\log n)^1\,\TS[n^{1466.92}] &&\text{Cor.~2.8, }x=1040000/709\\
11\quad&\subseteq(\exists\,n^0)^1(\forall\,n^{3156.82})^{3156.82}(\exists\,n^{2151.90})^{2151.90}(\forall\,n^{1466.85})^{1466.85}\,\BPTS[n^{2151.98}] &&\text{Lemma~5.1}\\
12\quad&\subseteq(\exists\,n^0)^1(\forall\,n^{3156.82})^{3156.82}(\exists\,n^{2151.90})^{2151.90}\,\BPTS[n^{3156.95}] &&\text{Lemma~5.1}\\
13\quad&\subseteq(\exists\,n^0)^1(\forall\,n^{3156.82})^{3156.82}\,\BPTS[n^{4631.25}] &&\text{Lemma~5.1}\\
14\quad&\subseteq(\exists\,n^0)^1\,\BPTS[n^{6794.04}] &&\text{Lemma~5.1}\\
15\quad&\subseteq\BPTS[n^{9966.86}] &&\text{Lemma~5.1}
\end{align*}
}

In lines 2, 3, 9 and 10, ``$\cdots$'' stands for the blocks of the previous line outside the ones shown. At line~8 the leading $n^0$ block of Theorem~2.10 is absorbed into $(\exists\,n^0)^1$, as the footnote to Corollary~2.11 prescribes; line~8 is the first use of rule~4. The final exponent is exactly $27410261545753261542320966055909/2750140100000000000000000000\approx9966.86<10^4$, so Lemma~5.2 applies.

A longer certificate at $c=8/5$ (2687 steps: $M(8)$ repeated 158 times, then the final slowdown) is also kernel-checked (\lean{MWL.mwcert\_ok}). The general theorem makes it redundant.

\section{The limit for \texorpdfstring{$c\ge\varphi$}{c >= phi}}\label{sec:limit}

\subsection{The relaxed calculus}

A \emph{relaxed line} (\lean{MW.MLine}) is $(s_1,\dots,s_m;\,k;\,a)$: the output exponents of the blocks (outermost first; Lean stores them innermost first), a kernel type $k\in\{\TS,\BPTS\}$ and a kernel exponent $a$. Quantifiers and guess exponents are dropped. The input exponent in front of the first block is $0$ (instead of $1$). The rules of \lean{MW.MStep}~$c$ only bound their outputs from below, and every speedup allows $0\le x\le a$. Write $s\cdot b$ for the list $s$ with $b$ appended (innermost), and $\pin'(s)$ for the last entry of $s$ ($0$ if $s$ is empty).
\begin{itemize}
\item \lean{rspeed0}: $(\emptyset;k;a)\vdash(b_1,b_2,b_3;\TS;a_2)$ for either $k$, with $b_1\ge0$, $b_2\ge\max(x,0)$, $b_3\ge0$, $a_2\ge a-x$.
\item \lean{rspeed}: $(s\cdot b;k;a)\vdash(s\cdot b\cdot b_2\cdot b_3;\TS;a_2)$ for either $k$, with $b_2\ge\max(b,x)$, $b_3\ge b$, $a_2\ge a-x$.
\item \lean{dspeed}: $(s\cdot b;\TS;a)\vdash(s\cdot t_2\cdot b_3;\TS;a_2)$ with $t_2\ge\max(b,x)$, $b_3\ge b$, $a_2\ge a-x$.
\item \lean{slow}: $(s\cdot b;k;a)\vdash(s;\BPTS;e)$ with $e\ge c\max(\pin'(s),b,a)$.
\end{itemize}
\lean{MW.UStep} adds \lean{rspeedU}: $(s\cdot b;k;a)\vdash(s\cdot b\cdot b_1\cdot b_2\cdot b_3;\TS;a_2)$ with $b_1\ge0$, $b_2\ge\max(b,x)$, $b_3\ge b$, $a_2\ge a-x$ (Theorem~2.10 under a block without absorbing its leading $n^0$ block). \lean{MW.StepW} adds weakening (\lean{MW.Weaker}): $(s;k;a)\vdash(s';k';a')$ with $s'$ of the same length as $s$ and $s'_i\ge s_i$, $a'\ge a$, and $k'=k$ or $k'=\BPTS$.

\begin{lemma}[\lean{MWL.lit\_m}, \lean{MWL.lit\_u}, \lean{MWL.lit\_w}]\label{lem:sim}
Let $c\ge0$. The encoding that keeps the output exponents and the kernel maps every \lean{MWL.Step} step between orderly lines with nonnegative exponents to an \lean{MW.MStep} step, every \lean{MWL.StepU} step to a \lean{MW.UStep} step and every \lean{MWL.StepW} step to a \lean{MW.StepW} step, and preserves these invariants.
\end{lemma}
\begin{proof}
Each relaxed rule has the literal outputs as admissible values. For the slowdown, the literal cost $c\max(a,o,g,\pin(\mathit{rest}))$ dominates $c\max(\pin'(\mathit{rest}),o,a)$ because $\pin(\emptyset)=1\ge0=\pin'(\emptyset)$. Guess exponents can be dropped because they satisfy $g\le o$ on every line.
\end{proof}

\subsection{The potential}

Let $c>0$ with $c^2\ge c+1$, and $g=1/(c^2-1)$. Then (\lean{MW.facts\_of\_phi})
\[ c\ge1,\qquad 0\le g\le1,\qquad 1+g\le c,\qquad 1+g=c^2g. \]
The third is equivalent to $(c-1)(c^2-1)\ge1$, i.e.\ $c(c^2-c-1)\ge0$, i.e.\ $c^2\ge c+1$; the fourth is an identity. Only the third condition constrains $c$: with this choice of $g$ the condition $1+g\le c^2g$, which the randomized speedup needs, holds with equality for every $c>1$. For a relaxed line define (\lean{MW.pot}, \lean{MW.up}, \lean{MW.}$\Phi$)
\begin{align*}
\pot(b,\BPTS,a)&=\max(b,a), & \pot(b,\TS,a)&=\max(b,\,g a),\\
\up(\emptyset;y)&=c\max(0,y), & \up(s\cdot b';y)&=\up(s;\,c\max(b',y)),\\
\Phi(\emptyset;k;a)&=a, & \Phi(s\cdot b;k;a)&=\up(s;\pot(b,k,a)).
\end{align*}
So $\Phi(s_1,\dots,s_m;k;a)=c\max(0,c\max(s_1,\dots c\max(s_{m-1},\pot(s_m,k,a))\cdots))$: the innermost block contributes a line value, and each outer block contributes the cost of a slowdown through it. $\up$ is monotone in $y$ and in every entry of the list for $c\ge0$ (\lean{MW.up\_mono}, \lean{MW.up\_mono\_list}), and $\pot$ is monotone in $b$ and $a$ for $g\ge0$ (\lean{MW.pot\_mono}). For $a\ge0$ and $g\le1$, $\pot(b,k,a)\le\max(b,a)$ (\lean{MW.pot\_le}).

\begin{lemma}[balance; \lean{MW.bal}, \lean{MW.key\_d}, \lean{MW.key\_r}]\label{lem:key}
For $0\le x\le a$ let $\mu=\max(x,\,g(a-x))$. Then $(1+g)\mu\ge g a$; hence $g a\le c\mu$ and $a\le c^2\mu$.
\end{lemma}
\begin{proof}
$(1+g)\mu=\mu+g\mu\ge g(a-x)+gx=ga$. Then $ga\le(1+g)\mu\le c\mu$; and $(1+g)a=c^2ga\le c^2(1+g)\mu$, so $a\le c^2\mu$.
\end{proof}

\begin{lemma}[the potential never decreases; \lean{MW.phi\_step}, \lean{MW.phi\_ustep}, \lean{MW.phi\_mono}]\label{lem:mono}
Let $c^2\ge c+1$, and let $L\vdash L'$ be a \lean{MW.StepW} step from a line with nonnegative entries. Then $\Phi(L)\le\Phi(L')$.
\end{lemma}
\begin{proof}
In each case the outer blocks $s$ are unchanged, so by monotonicity of $\up$ it suffices to compare the values passed to $\up(s;\cdot)$.
\begin{itemize}
\item \lean{dspeed} from $(s\cdot b;\TS;a)$: we need $\max(b,ga)\le c\max(t_2,b_3,ga_2)$. Now $b\le cb\le ct_2$, and $ga\le c\mu\le c\max(t_2,ga_2)$ by Lemma~\ref{lem:key}, since $t_2\ge x$ and $a_2\ge a-x$.
\item \lean{rspeed} from $(s\cdot b;k;a)$: since $\pot(b,k,a)\le\max(b,a)$, we need
\[ \max(b,a)\le c\max\bigl(b,\ c\max(b_2,b_3,ga_2)\bigr). \]
Now $b\le cb$, and $a\le c^2\mu\le c\cdot c\max(b_2,ga_2)$ by Lemma~\ref{lem:key}.
\item \lean{rspeedU}: the same with one more lift $c\max(b_1,\cdot)$ inserted, which only increases the value because $c\ge1$.
\item \lean{slow} from $(s\cdot b'\cdot b;k;a)$ to $(s\cdot b';\BPTS;e)$: before the step the value is $c\max(b',\pot(b,k,a))\le c\max(b',b,a)\le e$, and after it the value is $\pot(b',\BPTS,e)=\max(b',e)\ge e$. If $s\cdot b'$ is empty, $\Phi(L)=c\max(0,\pot(b,k,a))\le c\max(0,b,a)\le e=\Phi(L')$.
\item \lean{rspeed0} from a quantifier-free line: $\Phi(L')\ge c\cdot c\cdot c\mu\ge c\,a=c\,\Phi(L)\ge\Phi(L)$ (\lean{MW.first\_step}).
\item weakening: monotonicity of $\up$ and $\pot$, together with $\pot(b,\TS,a)\le\pot(b,\BPTS,a)$ for $a\ge0$, $g\le1$.\qedhere
\end{itemize}
\end{proof}

\begin{proof}[Proof of Theorem~\ref{thm:limit}]
Nonnegativity of all entries is preserved (\lean{MW.nn\_step}, \lean{MW.nn\_stepw}). A path from $\BPTS[n^{a_0}]$ stays at quantifier-free lines, with exponent at least $a_0$, while it uses only weakenings (\lean{MW.ustep\_nil\_ne}: a rule step from a quantifier-free line opens blocks). Its first rule step is therefore \lean{rspeed0} from a quantifier-free line of exponent $a\ge a_0$, and gives $\Phi\ge c\,a\ge c\,a_0$ (\lean{MW.first\_step}). By Lemma~\ref{lem:mono}, $\Phi$ never decreases afterwards, and on the final quantifier-free line $\Phi$ is its exponent $a'$. So $a'\ge ca_0>a_0$ when $a_0>0$, since $c>1$ (\lean{MW.qf\_ge}, \lean{MW.qf\_ge\_u}, \lean{MW.qf\_ge\_w}). The literal statements follow by Lemma~\ref{lem:sim}.
\end{proof}

\paragraph{Tightness.}
At $c=\varphi$, $g=1/(\varphi^2-1)=1/\varphi$ and $1+g=\varphi=c$, so the inequality for \lean{dspeed} is an equality when $x=g(a-x)$, and the one for \lean{rspeed} is then an equality as well. The macros of Section~\ref{sec:construction} approach this balance as their length grows: in $M(j+1)$ consecutive speedup parameters have ratio $c$, and at $c=\varphi$ the ratio $x/(a-x)$ of each deterministic speedup tends to $g=1/\varphi$. We note, as a heuristic and not as something the Lean proofs show, that a deterministic speedup can follow only a speedup in this calculus (every slowdown leaves a $\BPTS$ kernel), which is why the threshold is $\varphi$ and not the larger $2\cos(\pi/7)$ of the deterministic calculus \cite{BW15}.

\section{What the relaxed calculus covers}\label{sec:coverage}

Every relaxation below makes more steps possible, so the limit theorem covers it.
\begin{enumerate}
\item \emph{Unmerged randomized speedup.} Theorem~2.10 applied under a block $(Q\,n^a)^b$ whose leading $n^0$ block is not absorbed: \lean{rspeedU}. Since quantifiers are not tracked, this covers the leading block of either type; the literal \lean{MWL.StepU.rspeedU} formalises the alternating case $(Q\,n^a)^b(\neg Q\,n^0)^b(Q\,n^x)^{\max(b,x)}(\neg Q\,\log n)^b$.
\item \emph{Rule~4 with any quantifiers}, i.e.\ the printed Corollary~2.11 with $Q_{k+1},Q_{k+2}$ read as arbitrary, and \emph{rule~4 on a $\TS$ kernel} (through $\TS\subseteq\BPTS$): \lean{rspeed} for either kernel type.
\item \emph{Lemma~2.7 with $Q=\neg Q_k$ under a block}, the deterministic speedup that opens a new block of the opposite type, which \cite{MW21} set aside after Corollary~2.8 by \cite[Lemma~3.2]{W13}: it has the shape of \lean{rspeed} on a $\TS$ kernel.
\item \emph{Speedup parameters at the end points} $x=0$ and $x=a$, any output exponents above the exact values, and any output exponent of rule~3's blocks.
\item \emph{Input exponent $0$} instead of $1$ in front of the first block.
\item \emph{Weakening}: raising any exponent of an existing line, and $\TS[n^a]\subseteq\BPTS[n^a]$ (\lean{MW.StepW}; Lemma~\ref{lem:mono} proves that $\Phi$ is monotone under it). The proof of \cite[Theorem~5.3]{MW21} uses such steps.
\item \emph{Non-strict refutations}, $a'\le a_0$, as in the wording of Definition~2.14.
\end{enumerate}
Two things are outside both the formalisation and, as far as we can see, the source's notion of proof:
\begin{itemize}
\item \emph{Combining two adjacent blocks of the same type} as a separate step. Under the rules of \cite{MW21} lines always alternate; non-alternating lines arise only from the same-type variant of item~1. Applied immediately, combining turns that variant into the merged rule~4; combining at a later time is not formalised, and $\Phi$ is not monotone under it.
\item \emph{Contradictions other than $\BPTS[n^d]\subseteq\dots\subseteq\BPTS[n^{d'}]$}, for example a line with quantifiers contained in a smaller line of the same shape via an alternating hierarchy theorem. Definition~2.14 and Lemma~5.2 of \cite{MW21} use only the $\BPTS\to\BPTS$ form (Definition~2.14 also lists $\TS\to\TS$, which cannot occur in a proof whose slowdowns all produce $\BPTS$ kernels). Since $\Phi$ is homogeneous of degree~1 and never decreases, a path from a line $L$ to $\lambda L$ with $\lambda<1$ is impossible for $c\ge\varphi$; this is a paper argument, not checked in Lean, and the case of a target line that is smaller only in some coordinates is not analysed.
\end{itemize}
Rules at positions other than the innermost block are not rules of \cite{MW21} and are not included.

\section{Consequence for lower bounds}\label{sec:consequence}

The proof of \cite[Theorem~5.3]{MW21} uses: Theorem~2.10 (quoted from \cite{DvM06}, with a proof sketch), Corollaries~2.8 and~2.11, Lemma~5.1, and Lemma~5.2.

\begin{corollary}[conditional]\label{cor:lb}
Assume that these lemmas of \cite{MW21} hold as stated, and that each step of \lean{MWL.Step} is an instance of them (the modelling claim of Section~\ref{sec:calculus}). Then $\exists\cdot\BPTIME[n]\not\subseteq\BPTS[n^c]$ for every $c<\varphi$, and hence $\MATIME[n]\not\subseteq\BPTS[n^c]$.
\end{corollary}
\begin{proof}[Argument (not machine-checked)]
For $1\le c<\varphi$ Theorem~\ref{thm:construction} gives a refutation; for $c<1$ use the refutation at $c=1$, since $H_c$ implies $H_1$. Read as a chain of inclusions under $H_c$, each step is an instance of Theorem~2.10/Corollary~2.11, Corollary~2.8 or Lemma~5.1, and Lemma~5.2 turns the chain into a contradiction. Three points specific to the construction: (i) Corollary~2.11 is applied repeatedly under the outer block $(\exists\,n^0)^1$, each time absorbing a new $n^{o(1)}$-bit guess into it; the number of applications is a constant for fixed $c$, so the block stays of size $n^{o(1)}$. \cite{MW21} iterate under the same kind of block in the proof of their Theorem~5.4. (ii) The proof of Lemma~5.2 replaces the guess exponent $0$ of the first block by a small $\varepsilon>0$; in the construction that exponent enters only the final slowdown, whose cost $c\max(e',1,\varepsilon,1)$ is unchanged for $\varepsilon<1$. (iii) Lemma~5.2's proof sets $x'=x-(d-d')/2$, which must be positive; if the refutation ends far below $d$ (as the 8/5 certificate does), first weaken the last line to $\BPTS[n^{d''}]$ with $\max(d',d-2x)<d''<d$ (Section~\ref{sec:errata}).
\end{proof}

By $\MATIME[t]\subseteq\QCMATIME[t]$, which \cite[\S1.3]{MW21} state, the same conditional bound holds for $\QCMATIME[n]$; \cite[Main Theorem~3]{MW21} gives $1.5$ for that class against $\BPTS$.

The published bound $r_1$ depends on Theorem~2.10 and Lemma~5.2 in the same way as Corollary~\ref{cor:lb}. Lemma~5.2 rests on ``a fine-grained version of the alternating time hierarchy theorem (proved by direct diagonalization)'' for $n^{o(1)}$-space verifiers and quantifier blocks of $n^0$ and $x\log n$ type, which \cite{MW21} do not prove in detail; the ITCS version defers the proof of Lemma~5.2 to the full version. If Lemma~5.2 needed further hypotheses, both bounds would have to be re-examined.

\section{Two points in the source paper}\label{sec:errata}

We report two points in \cite{MW21} that we checked against the arXiv text. Neither affects the statements of \cite{MW21} or the results here.

\paragraph{The parameters in the proof of Theorem~5.3.}
The proof sets $x_1=d/c^3$ and $x_i=(c(1-\varepsilon))^{1-i}x_1$ with $\varepsilon=1/k$. With these values the displayed chain does not hold as written. Consecutive parameters have ratio $x_{i-1}/x_i=c(1-\varepsilon)<c$. The chain removes the block of guess $x_i$ from a kernel $cx_i$ and then uses $\BPTS[n^{c^2x_i}]\subseteq\BPTS[n^{cx_{i-1}}]$; together these need $cx_i\le x_{i-1}$, which fails. Moreover, if $c^2(c-1)<1$ (the range of the theorem), then $\sum_ix_i>d$ for all large $k$, so the kernel exponent $d-\sum_ix_i$ of the first line is negative. With $x_i=c^{k-i}T$ and $d=\sigma_kT$ (the macro $M(k)$ of Lemma~\ref{lem:mac}) the chain holds and gives exactly the stated condition in the limit. The script \texttt{scripts/mw\_check.py thm53} checks both claims exactly: with the printed parameters, $cx_i\le x_{i-1}$ fails for every $k$ tried and the kernel is negative for $c\in\{1.40,1.44,1.46\}$ and $k\ge5$; with the corrected ones, the best ratio of final to initial exponent over $k\le60$ is $0.998$ at $c=1.465$ and $1.0015$ at $c=1.466$.

\paragraph{The case $x'\le0$ in the proof of Lemma~5.2.}
The proof of Lemma~5.2 applies a speedup with parameter $x'=x-(d-d')/2$, where $x$ is the parameter of the first speedup. If $d-d'\ge2x$, then $x'\le0$ and the speedup is not defined. Weakening the end line first, $\BPTS[n^{d'}]\subseteq\BPTS[n^{d''}]$ with $\max(d',d-2x)<d''<d$, avoids this.

\section{What is and is not verified}\label{sec:verified}

\paragraph{(a) Kernel-checked.}
Lean~4.31 with Mathlib (pinned revision in the repository). The library builds with no \texttt{sorry}, no \texttt{admit}, no \texttt{native\_decide} and no new axioms; \texttt{\#print axioms} reports only \texttt{propext}, \texttt{Classical.choice} and \texttt{Quot.sound} for every headline theorem. The two concrete certificates are checked with \texttt{decide +kernel}, which adds no axiom. Checked are: the threshold of the \emph{formal} calculus in both directions (Theorems~\ref{thm:threshold}, \ref{thm:construction}, \ref{thm:limit}, including the variants with the unmerged randomized speedup and with weakening); the simulation of the literal rules in the relaxed calculus (Lemma~\ref{lem:sim}); monotonicity of the potential (Lemma~\ref{lem:mono}); soundness of the executable checker; the 15-step refutation at $c=1.467$ and the bound $r_1<1.467$; and the 2687-step refutation at $c=8/5$.

\paragraph{(b) Not machine-checked.}
\begin{itemize}
\item \emph{That the formal rules coincide with the rules of \cite{MW21}.} This is a modelling claim, compared rule by rule in Table~\ref{tab:rules}, with the reading points listed after it. The construction uses only \lean{MWL.Step}, whose rules are the exact rules with equal or stricter side conditions. The limit side is proved for a relaxation (Section~\ref{sec:coverage}).
\item \emph{Soundness of the lemmas of \cite{MW21}}: Theorem~2.10 (labelled a sketch there, quoting \cite{DvM06}), Corollaries~2.8 and~2.11, Lemma~5.1 and Lemma~5.2. The published bound $r_1$ depends on them in the same way; in particular Lemma~5.2 is only sketched (Section~\ref{sec:consequence}). We do not formalise complexity classes.
\item \emph{Corollary~\ref{cor:lb}} and the points (i)--(iii) in its argument.
\item The paper arguments marked as such in Section~\ref{sec:coverage} (contradictions of other forms) and the heuristic remark after the limit proof.
\item The scripts in \texttt{scripts/} use exact rational arithmetic and are independent of the Lean code, but they are not part of the kernel-checked results.
\end{itemize}

\paragraph{(c) Scope.}
The result is a limit on one proof calculus, the one of \cite[Definition~2.14 and \S5]{MW21}, with the variants listed in Section~\ref{sec:coverage}. Other speedup or slowdown rules (for instance Grover search, as in \cite[\S4]{MW21}), other contradiction forms, and other proof methods are not covered. Nothing here bears on lower bounds for $\exists\cdot\BPTIME$ or $\MATIME$ proved by other methods.

\section{Relation to earlier work}

For SAT against deterministic $n^{o(1)}$-space algorithms, Fortnow and van Melkebeek \cite{FvM00,FLvMV05} obtained $\varphi$; Williams later reached $2\cos(\pi/7)$, which Buss and Williams \cite{BW15} proved optimal for alternation-trading proofs with those rules. For $\MATIME$ against randomized $n^{o(1)}$-space algorithms, Diehl and van Melkebeek \cite{DvM06} and Diehl \cite{Die07} obtained $\sqrt2$, and Mudigonda and Williams \cite{MW21} obtained $r_1\approx1.4656$. Iterating the macro $M(1)$ reproduces $\sqrt2$, using $M(k)$ once reproduces $r_1$ as $k\to\infty$, and iterating $M(k)$ gives every $c<\varphi$; $\varphi$ is the exact threshold of the calculus.

\section*{Licence}
This text is licensed under CC BY 4.0. The Lean code and the scripts are licensed under the Apache License~2.0.

\begin{thebibliography}{FLvMV05}

\bibitem[BW15]{BW15}
S.~R. Buss and R.~Williams.
\newblock Limits on alternation trading proofs for time--space lower bounds.
\newblock \emph{Computational Complexity}, 24(3):533--600, 2015.
\newblock \href{https://doi.org/10.1007/s00037-015-0104-9}{doi:10.1007/s00037-015-0104-9}.

\bibitem[Die07]{Die07}
S.~Diehl.
\newblock Lower bounds for swapping Arthur and Merlin.
\newblock In \emph{Approximation, Randomization, and Combinatorial Optimization (APPROX-RANDOM 2007)}, \emph{Lecture Notes in Computer Science}, pages 449--463. Springer, 2007.

\bibitem[DvM06]{DvM06}
S.~Diehl and D.~van Melkebeek.
\newblock Time-space lower bounds for the polynomial-time hierarchy on randomized machines.
\newblock \emph{SIAM Journal on Computing}, 36(3):563--594, 2006.
\newblock \href{https://doi.org/10.1137/050642228}{doi:10.1137/050642228}.

\bibitem[FLvMV05]{FLvMV05}
L.~Fortnow, R.~Lipton, D.~van Melkebeek, and A.~Viglas.
\newblock Time-space lower bounds for satisfiability.
\newblock \emph{Journal of the ACM}, 52(6):835--865, 2005.
\newblock \href{https://doi.org/10.1145/1101821.1101822}{doi:10.1145/1101821.1101822}.

\bibitem[FvM00]{FvM00}
L.~Fortnow and D.~van Melkebeek.
\newblock Time-space tradeoffs for nondeterministic computation.
\newblock In \emph{Proceedings of the 15th Annual IEEE Conference on Computational Complexity (CCC 2000)}, pages 2--13, 2000.
\newblock \href{https://doi.org/10.1109/CCC.2000.856730}{doi:10.1109/CCC.2000.856730}.

\bibitem[MW21]{MW21}
A.~S. Mudigonda and R.~R. Williams.
\newblock Time-space lower bounds for simulating proof systems with quantum and randomized verifiers.
\newblock In \emph{12th Innovations in Theoretical Computer Science Conference (ITCS 2021)}, volume 185 of \emph{LIPIcs}, pages 50:1--50:20, 2021.
\newblock \href{https://doi.org/10.4230/LIPIcs.ITCS.2021.50}{doi:10.4230/LIPIcs.ITCS.2021.50}.
\newblock Full version: \href{https://arxiv.org/abs/2012.00330}{arXiv:2012.00330} (v2, January 2021), whose numbering we use.

\bibitem[W13]{W13}
R.~Williams.
\newblock Alternation-trading proofs, linear programming, and lower bounds.
\newblock \emph{ACM Transactions on Computation Theory}, 5(2):Article~6, 1--49, 2013.
\newblock \href{https://doi.org/10.1145/2493246.2493249}{doi:10.1145/2493246.2493249}.

\end{thebibliography}

\end{document}
